\chapter{הפקת מסמכים מקצועיים באמצעות LaTeX}
\section{מ-artifacts מאומתים ל-LaTeX: templates ו-escaping}
לאחר שכל ה-Agents סיימו את עבודתם ושלב ה-validation אישר את התוצרים, עומד לרשותנו אוסף של artifacts מובנים: כתב היד שעבר review, מטא-נתונים, נכסים גרפיים וקובץ מקורות. השלב הבא אינו שלב יצירתי אלא דטרמיניסטי לחלוטין, ובכך טמון עיקרון מרכזי בתכנון המערכת. ה-Agent מתאר כוונה בלבד, באמצעות מסמך latex\_spec מובנה הקובע איזה template לבחור, כיצד למפות פרקים וכיצד לשבץ נכסים, אך לעולם אינו כותב LaTeX גולמי בעצמו. את הרינדור מבצע רכיב Python ייעודי, ה-renderer, הממלא templates של Jinja2 בתוכן המאומת ומפיק את הקובץ main.tex ואת קבצי הפרקים. הפרדה זו בין שיקול הדעת של ה-Agent לבין ההפקה הדטרמיניסטית היא לב ההבטחה לשחזוריות (reproducibility). אתגר קריטי בשלב זה הוא escaping: טקסט שמקורו ב-Agent נחשב untrusted text, ותווים מיוחדים של LaTeX, כדוגמת סימני אחוז, אמפרסנד או סוגריים מסולסלים, עלולים לשבור את הקומפילציה ואף לאפשר הזרקת פקודות. לכן ה-renderer מחזיק מודול escaping ייעודי, המעבד כל מחרוזת לפני שיבוצה ב-template ומבטיח שתוכן שרירותי לא יהפוך לקוד פעיל. כך מתקבל קוד מקור תקני, אחיד וצפוי, שאינו תלוי בגחמות הניסוח של המודל. גישה זו תומכת גם ב-observability: מאחר שה-renderer ידוע ומבוקר, ניתן להריץ אותו אף ללא toolchain מותקן ולקבל קובצי tex תקפים לבדיקה. כך מופרד בבירור שלב ההרכבה משלב הקומפילציה בפועל, וכל סטייה בתוכן נחשפת עוד לפני ההידור.

\section{כתיבה דו-כיוונית עברית-אנגלית עם מנוע Unicode}
מסמך מקצועי בעברית המשלב מונחים טכניים באנגלית מציב אתגר טיפוגרפי של ממש, שכן עברית נכתבת מימין לשמאל ואילו מונחים כדוגמת Agent, Crew או validation נכתבים משמאל לימין באותה שורה. ניהול נכון של מצב זה, המכונה BiDi או bidirectional typesetting, מחייב מנוע מודרני המבוסס על Unicode. מנוע pdfLaTeX הוותיק אינו מספק תמיכה אמינה בכך, ולכן המערכת בוחרת ב-LuaLaTeX כברירת מחדל, עם XeLaTeX כ-fallback. שני מנועים אלה קוראים ישירות קבצים בקידוד Unicode ויודעים לטעון גופנים מודרניים התומכים בעברית, וכך מתייתרות שכבות הקידוד המסורבלות של העבר. את ניהול השפות מבצעת חבילת polyglossia יחד עם רכיב bidi, המחליפים את החבילות הישנות שלא הותאמו לעברית. polyglossia מאפשרת להגדיר שפה ראשית ושפה משנית ולעבור ביניהן באופן מפורש, כך שכל קטע טקסט מקבל את כיוון הכתיבה הנכון לו. העיקרון המעשי הוא בידוד: יש לעטוף כל רצף באנגלית בתוך משפט עברי כך שיישמר כיחידה רציפה משמאל לימין, ולמנוע מהאלגוריתם הדו-כיווני לערבב את סדר האותיות סביב נקודות המעבר. בקרת איכות מחייבת אימות חזותי של המעברים: יש לוודא שמספרים, סימני פיסוק וסוגריים נצמדים לצד הנכון, ושמונחים טכניים אינם נשברים. ה-renderer מקבל מטא-נתוני BiDi מתוך ה-latex\_spec ומשבץ את מתגי השפה במקומותיהם הנכונים, כך שגם תוצר רב-לשוני זה נשאר דטרמיניסטי וניתן לשחזור מלא.

\section{ביבליוגרפיה וקומפילציה רב-מעברית של citations}
אחד ההיבטים הפחות אינטואיטיביים בהפקת מסמך LaTeX הוא שאי אפשר להפיק ביבליוגרפיה תקינה במעבר קומפילציה יחיד. הסיבה נעוצה באופן שבו LaTeX פותר הפניות: במעבר הראשון המהדר אינו יודע עדיין לאילו מקורות מפנה הטקסט, ולכן הוא רק אוסף את סימני ה-citation ורושם אותם בקובצי עזר, כדוגמת קובץ aux וקובץ bcf. רק לאחר מכן נכנס לפעולה ה-backend הביבליוגרפי, biber, הקורא את רשימת ההפניות שנאספה, מצליב אותה מול קובץ references.bib, וממיין ומעצב את הרשומות לפי סגנון הציטוט הנדרש. המעברים הבאים של המהדר משלבים פלט זה חזרה אל המסמך, וממלאים את מספרי ההפניות ואת רשימת המקורות במקומותיהם. לכן הרצף הקנוני במערכת הוא ריצת המהדר, אחריה biber, ואחריה שתי ריצות נוספות של המהדר. ללא רצף זה יישארו במסמך סימני שאלה כפולים במקום מספרי citation, וגם cross-references פנימיים, כדוגמת הפניות לאיורים ולטבלאות, לא ייפתרו. כאן בולט שוב עקרון ההפרדה: ניהול ה-citations עצמו דטרמיניסטי ואינו Agent, שכן אסור שמודל ימציא מפתחות או רשומות. קובץ ה-bib נוצר מתוך רישום מקורות מבוקר, ושלב reconciliation מאמת שכל citation marker במסמך מתאים למקור מוכר, ונכשל בקול רם אם נמצא מפתח שאינו קיים. כך משלבת המערכת קומפילציה רב-מעברית עם validation מחמיר, ומבטיחה מסמך אקדמי אמין שבו כל הפניה נפתרת כראוי.

לקריאה נוספת בנושאי פרק זה, ראו \cite{langchain_docs}.

